\documentclass[8pt,a4paper,landscape]{extarticle}

\usepackage[utf8]{inputenc}
\usepackage[T1]{fontenc}
\usepackage[brazil,shorthands=off]{babel}
\usepackage{lmodern}
\usepackage[scaled=0.82]{beramono}
\IfFileExists{upquote.sty}{\usepackage{upquote}}{}

\usepackage[
  a4paper,
  landscape,
  left=0.55cm,
  right=0.55cm,
  top=1.05cm,
  bottom=0.55cm,
  headheight=11pt,
  headsep=0.22cm
]{geometry}

\usepackage{multicol}
\usepackage{amsmath,amssymb,amsfonts,amsthm,mathtools}
\usepackage{xcolor}
\usepackage{listings}
\usepackage{fancyhdr}
\usepackage{titlesec}
\usepackage{enumitem}
\usepackage{booktabs}
\usepackage{array}
\usepackage{microtype}
\usepackage[hidelinks,bookmarks=false,pdfencoding=auto]{hyperref}

% Paleta de cores limpa estilo KACTL
\definecolor{dkgreen}{rgb}{0.0,0.40,0.16}
\definecolor{dkblue}{rgb}{0.04,0.18,0.62}
\definecolor{dkred}{rgb}{0.58,0.06,0.06}
\definecolor{dkpurple}{rgb}{0.42,0.06,0.48}
\definecolor{codebg}{rgb}{0.985,0.988,0.992}
\definecolor{boxbg}{rgb}{0.945,0.955,0.972}
\definecolor{rulecolor}{rgb}{0.72,0.75,0.80}
\definecolor{secblue}{rgb}{0.06,0.20,0.42}

\lstdefinestyle{icpccpp}{
  language=C++,
  basicstyle=\ttfamily\fontsize{6.2pt}{7.0pt}\selectfont,
  keywordstyle=\color{dkblue}\bfseries,
  keywordstyle=[2]\color{dkpurple}\bfseries,
  commentstyle=\color{dkgreen}\itshape,
  stringstyle=\color{dkred},
  numbers=none,
  backgroundcolor=\color{codebg},
  frame=single,
  rulecolor=\color{rulecolor},
  framesep=1.5pt,
  xleftmargin=2pt,
  xrightmargin=2pt,
  aboveskip=2pt,
  belowskip=2pt,
  tabsize=2,
  breaklines=true,
  breakatwhitespace=false,
  showstringspaces=false,
  keepspaces=true,
  columns=fixed,
  basewidth=0.50em,
  morekeywords={
    constexpr,decltype,noexcept,nullptr,static_assert,
    concept,requires,__int128,__int128_t,__uint128_t,
    ll,ull,i128,u128,ld,pii,pll,vi,vll
  },
  morekeywords=[2]{
    rep,all,sz,pb,eb,fst,snd,chmin,chmax
  },
  literate=
    {á}{{\'a}}1 {é}{{\'e}}1 {í}{{\'i}}1 {ó}{{\'o}}1 {ú}{{\'u}}1
    {Á}{{\'A}}1 {É}{{\'E}}1 {Í}{{\'I}}1 {Ó}{{\'O}}1 {Ú}{{\'U}}1
    {à}{{\`a}}1 {ã}{{\~a}}1 {õ}{{\~o}}1 {ç}{{\c{c}}}1 {Ç}{{\c{C}}}1
    {â}{{\^a}}1 {ê}{{\^e}}1 {ô}{{\^o}}1
}

\lstdefinestyle{icpcplain}{
  basicstyle=\ttfamily\fontsize{6.2pt}{7.0pt}\selectfont,
  backgroundcolor=\color{codebg},
  frame=single,
  rulecolor=\color{rulecolor},
  framesep=1.5pt,
  xleftmargin=2pt,
  xrightmargin=2pt,
  aboveskip=2pt,
  belowskip=2pt,
  tabsize=4,
  breaklines=true,
  showstringspaces=false,
  keepspaces=true,
  columns=fixed,
  basewidth=0.50em
}

\lstset{style=icpccpp}

\setlength{\columnsep}{0.42cm}
\setlength{\columnseprule}{0.3pt}
\def\columnseprulecolor{\color{rulecolor}}
\setlength{\parindent}{0pt}
\setlength{\parskip}{1.2pt}

\setlist[itemize]{leftmargin=9pt,itemsep=0.4pt,topsep=1pt,parsep=0pt}
\setlist[enumerate]{leftmargin=11pt,itemsep=0.4pt,topsep=1pt,parsep=0pt}

\titleformat{\section}
  {\normalfont\fontsize{9.0pt}{10.5pt}\selectfont\bfseries\color{secblue}}
  {\thesection.}{0.35em}{}[\vspace{-1pt}\titlerule]
\titlespacing*{\section}{0pt}{3.5pt}{1.8pt}

\titleformat{\subsection}
  {\normalfont\fontsize{7.8pt}{9.0pt}\selectfont\bfseries\color{black}}
  {\thesubsection}{0.35em}{}
\titlespacing*{\subsection}{0pt}{2.8pt}{1pt}

\titleformat{\subsubsection}
  {\normalfont\fontsize{7.2pt}{8.2pt}\selectfont\bfseries\color{secblue!85!black}}
  {\thesubsubsection}{0.3em}{}
\titlespacing*{\subsubsection}{0pt}{2pt}{0.5pt}

% Cabecalho sem colisao: Esquerda = Titulo curto | Direita = Secao atual + Pagina
\pagestyle{fancy}
\fancyhf{}
\renewcommand{\sectionmark}[1]{\markboth{\thesection.\ #1}{}}
\fancyhead[L]{\footnotesize\textbf{UFABC — Team Reference Document (TRD)} $\mid$ \textit{ICPC 2026}}
\fancyhead[R]{\footnotesize\textit{\nouppercase{\leftmark}} \quad \textbf{Pág.\ \thepage}}
\renewcommand{\headrulewidth}{0.4pt}
\renewcommand{\footrulewidth}{0pt}

\newcommand{\Ocomp}[1]{\mathcal{O}(#1)}
\newcommand{\eps}{\varepsilon}
\newcommand{\MOD}{10^9+7}
\newcommand{\NTTMOD}{998244353}
\newcommand{\ZZ}{\mathbb{Z}}
\newcommand{\RR}{\mathbb{R}}

% Caixa teorica enxuta com \raggedright (elimina 100% dos Underfull \hbox)
\newcommand{\algobox}[4]{%
  \vspace{1pt}%
  \noindent\colorbox{boxbg}{%
    \parbox{\dimexpr\linewidth-2\fboxsep\relax}{%
      \raggedright\fontsize{6.5pt}{7.6pt}\selectfont%
      \textbf{\color{secblue}Uso:} #1 \par%
      \textbf{\color{secblue}Ideia:} #2 \par%
      \textbf{\color{dkred}Armadilhas:} #3 $\mid$ \textbf{\color{secblue}Complexidade:} #4%
    }%
  }%
  \vspace{1pt}\par%
}

\begin{document}
\begin{multicols*}{3}

\begin{center}
  {\large\bfseries\color{secblue} UFABC — ICPC Team Reference Document}\\[1pt]
  {\scriptsize\textbf{C++23 \& Python 3 Competitive Programming Handbook — Fase Nacional 2026}}
\end{center}
\vspace{-3pt}
{\scriptsize\tableofcontents}
\vspace{2pt}
\hrule
\vspace{2pt}

% ============================================================================
% SECAO 0: AMBIENTE, TEMPLATES, MAKEFILE POSIX, COMPILACAO E EDITORES
% ============================================================================
\section*{0. Ambiente, Templates, Makefile e Editores}
\addcontentsline{toc}{section}{0. Ambiente, Templates, Makefile e Editores}
\markboth{0.\ Ambiente, Templates, Makefile e Editores}{}

\subsection*{0.1 Templates Oficiais (C++23 \& Python 3)}

\textbf{Template C++ (\texttt{\textasciitilde/.template.cpp}):}
\begin{lstlisting}
#include <bits/stdc++.h>
typedef long long ll;
using namespace std;

void task() {

}

int main() {
    std::ios_base::sync_with_stdio(false);
    std::cin.tie(NULL);

    ll t = 1;
    // cin >> t;
    while (t--) {
        task();
    }
    return 0;
}
\end{lstlisting}

\textbf{Template Python (\texttt{solution.py}):}
\begin{lstlisting}[style=icpcplain]
import sys
sys.setrecursionlimit(2**32)

def main():
    pass

main()
\end{lstlisting}

\subsection*{0.2 Automação com Makefile POSIX}
\algobox
  {Automação portável compatível com GNU Make (\texttt{gmake}), BSD Make (\texttt{bmake}), BusyBox Make e POSIX Make.}
  {\texttt{.POSIX:} força conformidade estrita via \texttt{/bin/sh}. \texttt{.SUFFIXES:} vazio limpa regras legadas e registra apenas \texttt{.cpp} e \texttt{.debug}. \texttt{[A-Z]} no \texttt{clean} apaga só os binários das questões.}
  {Usar caractere \texttt{TAB} real na indentação das regras. \texttt{make A} gera binário rápido (\texttt{-O3}); \texttt{make A.debug} gera binário com Sanitizers.}
  {Uso: \texttt{make A \&\& ./A < in}}

\begin{lstlisting}[style=icpcplain]
.POSIX:

CXX      = c++
CXXFLAGS = -O3 -Wall -Wextra -Wconversion -Wshadow -std=c++23
DBGFLAGS = -fsanitize=address,undefined -g -Wall -Wshadow -std=c++23

.SUFFIXES:
.SUFFIXES: .cpp .debug

.cpp:
	$(CXX) $(CXXFLAGS) $< -o $@

.cpp.debug:
	$(CXX) $(DBGFLAGS) $< -o $@

clean:
	rm -f [A-Z] *.debug
\end{lstlisting}

\subsection*{0.3 Guia de Compilação \& Execução}
\begin{itemize}
  \item \textbf{Flags Essenciais C++23:}
    \texttt{-O3} (otimização máxima), \texttt{-Wall -Wextra} (avisos e UB), \texttt{-Wconversion} (alerta conversão implícita como \texttt{ll} $\to$ \texttt{int}), \texttt{-Wshadow} (variável local mascarando global), \texttt{-std=c++23}.
  \item \textbf{Linux (GCC):}
\begin{lstlisting}[style=icpcplain]
g++ -O3 -Wall -Wextra -Wconversion -Wshadow -std=c++23 sol.cpp -o sol
c++ -O3 -Wall -Wextra -Wconversion -Wshadow -std=c++23 sol.cpp -o sol
\end{lstlisting}
  \item \textbf{FreeBSD (Clang):}
\begin{lstlisting}[style=icpcplain]
clang++ -O3 -Wall -Wextra -Wconversion -Wshadow -std=c++23 sol.cpp -o sol
c++ -O3 -Wall -Wextra -Wconversion -Wshadow -std=c++23 sol.cpp -o sol
CC  -O3 -Wall -Wextra -Wconversion -Wshadow -std=c++23 sol.cpp -o sol
\end{lstlisting}
  \item \textbf{Python 3 \& PyPy 3 (JIT — até 10x mais rápido contra TLE):}
\begin{lstlisting}[style=icpcplain]
python3 solution.py < in
pypy3   solution.py < in
\end{lstlisting}
  \item \textbf{Depuração Local com Sanitizers (Out-of-bounds / UB / Leak):}
\begin{lstlisting}[style=icpcplain]
c++ -fsanitize=address,undefined -g -Wall -Wshadow -std=c++23 sol.cpp -o debug
./debug < in
\end{lstlisting}
\end{itemize}

\subsection*{0.4 Configuração de Editores no Warmup}

\textbf{Vim (\texttt{\textasciitilde/.vimrc} — 4 linhas):}
\begin{lstlisting}[style=icpcplain]
syn on|se nu rnu cin noet ts=4 sw=4 sts=4 bs=2 is hls ww+=<,>,h,l,[,]
let &t_SI="\<Esc>[5 q"|let &t_EI="\<Esc>[1 q"
au BufNewFile *.cpp 0r ~/.template.cpp
map <F5> :w<CR>:!c++ -O3 -std=c++23 % -o %:r && ./%:r < in<CR>
\end{lstlisting}
\begin{itemize}
  \item \texttt{nu rnu}: linha absoluta na atual e relativa nas demais; \texttt{cin noet ts=4 sw=4 sts=4}: indentação C/C++ com tab de 4 colunas; \texttt{t\_SI/t\_EI}: cursor barra (\texttt{|}) no Insert e bloco no Normal.
  \item \textbf{Compilação manual no Vim:} \texttt{:!c++ -O3 -std=c++23 \% -o \%:r \&\& ./\%:r < in}
\end{itemize}

\textbf{GNU Emacs (\texttt{\textasciitilde/.emacs} ou \texttt{*scratch*} com \texttt{M-x eval-buffer}):}
\begin{lstlisting}[style=icpcplain]
(global-display-line-numbers-mode 1)
(setq-default indent-tabs-mode nil)
(setq-default c-basic-offset 4)
(setq compilation-scroll-output t)
(global-set-key (kbd "<f5>") 'compile)
\end{lstlisting}
\begin{itemize}
  \item \textbf{Fluxo \texttt{<F5>} (\texttt{M-x compile}):} \texttt{c++ -O3 -Wall -Wextra -std=c++23 sol.cpp -o sol \&\& ./sol < in}
  \item \textbf{Navegação de Erros:} \texttt{C-x `} ou \texttt{M-g n} (próximo erro), \texttt{M-g p} (erro anterior), \texttt{g} no buffer \texttt{*compilation*} (recompilar).
\end{itemize}

\subsection*{0.5 Macros dos Capítulos, \texttt{\_\_int128}, Hash e Bits}
\begin{lstlisting}
using ll = long long;
using ull = unsigned long long;
using i128 = __int128_t;
using u128 = __uint128_t;
using ld = long double;
using pii = pair<int, int>;
using pll = pair<ll, ll>;
using vi = vector<int>;
using vll = vector<ll>;

#define rep(i, a, b) for (int i = (a); i < (int)(b); ++i)
#define all(x) begin(x), end(x)
#define sz(x) (int)size(x)
#define pb push_back
#define eb emplace_back
#define fst first
#define snd second

template <class T> bool chmin(T &a, const T &b) { return b < a ? a = b, 1 : 0; }
template <class T> bool chmax(T &a, const T &b) { return a < b ? a = b, 1 : 0; }
inline i128 abs128(i128 x) { return x < 0 ? -x : x; }

ll floor_div(ll a, ll b) {
  if (b < 0) a = -a, b = -b;
  return a / b - (a % b < 0);
}
ll ceil_div(ll a, ll b) {
  if (b < 0) a = -a, b = -b;
  return a / b + (a % b > 0);
}

mt19937_64 rng(chrono::steady_clock::now().time_since_epoch().count());
struct CustomHash {
  static ull splitmix64(ull x) {
    x += 0x9e3779b97f4a7c15ULL;
    x = (x ^ (x >> 30)) * 0xbf58476d1ce4e5b9ULL;
    x = (x ^ (x >> 27)) * 0x94d049bb133111ebULL;
    return x ^ (x >> 31);
  }
  size_t operator()(ull x) const {
    static const ull FIXED_RANDOM = rng();
    return splitmix64(x + FIXED_RANDOM);
  }
  template <class A, class B>
  size_t operator()(const pair<A, B>& p) const {
    ull h1 = (*this)(p.first), h2 = (*this)(p.second);
    return h1 ^ (h2 + 0x9e3779b97f4a7c15ULL + (h1 << 6) + (h1 >> 2));
  }
};
\end{lstlisting}

% ============================================================================
% INCLUSAO MODULAR DOS CAPITULOS
% ============================================================================
\InputIfFileExists{cap1.tex}{}{}
\InputIfFileExists{cap2.tex}{}{}
\InputIfFileExists{cap3.tex}{}{}
\InputIfFileExists{cap4.tex}{}{}
\InputIfFileExists{cap5.tex}{}{}
\InputIfFileExists{cap6.tex}{}{}
\InputIfFileExists{cap7.tex}{}{}
\InputIfFileExists{cap8.tex}{}{}

\end{multicols*}
\end{document}
